% GENERATED FILE. Edit canonical lessons, then run npm run generate:books.
% canonical-source-hash: fnv1a64:ef5fc3fbf91d5284
% canonical-lessons: GE-C316-erzaehler, GE-C316-handlung, GE-C316-figur, GE-C316-kapitel, GE-C316-fortsetzung

\chapter{Wer erzählt?}
\label{ch:ge-wer-erzahlt}
\hlchaptermodality{316}

\begin{chapteropening}
\textbf{By the end of this chapter:} \emph{I can talk about how a story is built: a narrator, the plot, a character, a chapter and a sequel.}

Complete the last lesson of chapter 316: I can talk about how a story is built: a narrator, the plot, a character, a chapter and a sequel.
\end{chapteropening}

\section[der Erzähler]{der Erzähler — the narrator, the storyteller}

\label{lesson:GE-C316-erzaehler}



\begin{quote}
Before the new one: say the German for the giant, then the German for the dwarf.
\end{quote}

\subsection*{You'll want to know: der Erzähler}
\textbf{der Erzähler} — \textquotedblleft{}the narrator, the storyteller\textquotedblright{}.

\begin{quote}
Der Erzähler spricht sehr leise.
\end{quote}

\begin{quote}
\emph{The narrator speaks very quietly.}
\end{quote}

Plural: \emph{die Erzähler}. Feminine \emph{die Erzählerin}.

\begin{cousinweb}
\emph{Erzähler} is built on \emph{erzählen} \textquotedblleft{}to tell\textquotedblright{}, from \emph{zählen} \textquotedblleft{}to count\textquotedblright{}, Old High German \emph{zellen}, Proto-Germanic \emph{taljaną} (EtymDB): to count out events one by one.
\end{cousinweb}

\begin{grammarlens}[title={What you've built}]
One more word for telling, explaining and describing at B1.
\end{grammarlens}

\subsection*{Your turn}
\begin{itemize}
\raggedright
  \item \emph{Say it:} \emph{der Erzähler}
  \item \emph{Say it:} \emph{der Erzähler}, once more
  \item \emph{Say it:} a sentence of your own with \emph{der Erzähler}
  \item \emph{Recall:} say the German for the giant, then the German for the dwarf, then say \emph{der Erzähler} again
\end{itemize}

\begin{tcolorbox}[breakable,colback=teal!4,colframe=teal!35!black,title={Before you move on}]
What does der Erzähler mean? (\textquotedblleft{}The narrator, the storyteller\textquotedblright{}.) Say it once more.
\end{tcolorbox}
\section[die Handlung]{die Handlung — the plot, the action}

\label{lesson:GE-C316-handlung}



\begin{quote}
Before the new one: say the German for the dwarf, then the German for the narrator.
\end{quote}

\subsection*{You'll want to know: die Handlung}
\textbf{die Handlung} — \textquotedblleft{}the plot, the action\textquotedblright{}.

\begin{quote}
Die Handlung des Films war spannend.
\end{quote}

\begin{quote}
\emph{The plot of the film was exciting.}
\end{quote}

Plural: \emph{die Handlungen}. Also an act or deed.

\begin{cousinweb}
\emph{Handlung} is built on \emph{handeln} \textquotedblleft{}to act\textquotedblright{}, from Old High German \emph{hantalōn}, which EtymDB derives from Proto-Germanic \emph{handuz} \textquotedblleft{}hand\textquotedblright{}.
\end{cousinweb}

\begin{grammarlens}[title={What you've built}]
One more word for telling, explaining and describing at B1.
\end{grammarlens}

\subsection*{Your turn}
\begin{itemize}
\raggedright
  \item \emph{Say it:} \emph{die Handlung}
  \item \emph{Say it:} \emph{die Handlung}, once more
  \item \emph{Say it:} a sentence of your own with \emph{die Handlung}
  \item \emph{Recall:} say the German for the dwarf, then the German for the narrator, then say \emph{die Handlung} again
\end{itemize}

\begin{tcolorbox}[breakable,colback=teal!4,colframe=teal!35!black,title={Before you move on}]
What does die Handlung mean? (\textquotedblleft{}The plot, the action\textquotedblright{}.) Say it once more.
\end{tcolorbox}
\section[die Figur]{die Figur — the character, the figure}

\label{lesson:GE-C316-figur}



\begin{quote}
Before the new one: say the German for the narrator, then the German for the plot.
\end{quote}

\subsection*{You'll want to know: die Figur}
\textbf{die Figur} — \textquotedblleft{}the character, the figure\textquotedblright{}.

\begin{quote}
Die wichtigste Figur im Buch ist ein Junge.
\end{quote}

\begin{quote}
\emph{The most important character in the book is a boy.}
\end{quote}

Plural: \emph{die Figuren}. Also a person's figure, or a chess piece.

\begin{cousinweb}
\emph{Figur} was borrowed from Latin \emph{figūra} \textquotedblleft{}a shape, a form\textquotedblright{} (EtymDB), built on \emph{fingere} \textquotedblleft{}to form, to shape\textquotedblright{} (Lewis \& Short).
\end{cousinweb}

\begin{grammarlens}[title={What you've built}]
One more word for telling, explaining and describing at B1.
\end{grammarlens}

\subsection*{Your turn}
\begin{itemize}
\raggedright
  \item \emph{Say it:} \emph{die Figur}
  \item \emph{Say it:} \emph{die Figur}, once more
  \item \emph{Say it:} a sentence of your own with \emph{die Figur}
  \item \emph{Recall:} say the German for the narrator, then the German for the plot, then say \emph{die Figur} again
\end{itemize}

\begin{tcolorbox}[breakable,colback=teal!4,colframe=teal!35!black,title={Before you move on}]
What does die Figur mean? (\textquotedblleft{}The character, the figure\textquotedblright{}.) Say it once more.
\end{tcolorbox}
\section[das Kapitel]{das Kapitel — the chapter}

\label{lesson:GE-C316-kapitel}



\begin{quote}
Before the new one: say the German for the plot, then the German for the character.
\end{quote}

\subsection*{You'll want to know: das Kapitel}
\textbf{das Kapitel} — \textquotedblleft{}the chapter\textquotedblright{}.

\begin{quote}
Im letzten Kapitel findet er seinen Vater.
\end{quote}

\begin{quote}
\emph{In the last chapter he finds his father.}
\end{quote}

Plural: \emph{die Kapitel}.

\begin{cousinweb}
\emph{Kapitel} is a learned borrowing from Latin \emph{capitulum} \textquotedblleft{}a little head, a section\textquotedblright{}, from \emph{caput} \textquotedblleft{}head\textquotedblright{} (Lewis \& Short); EtymDB lists no parent.
\end{cousinweb}

\begin{grammarlens}[title={What you've built}]
One more word for telling, explaining and describing at B1.
\end{grammarlens}

\subsection*{Your turn}
\begin{itemize}
\raggedright
  \item \emph{Say it:} \emph{das Kapitel}
  \item \emph{Say it:} \emph{das Kapitel}, once more
  \item \emph{Say it:} a sentence of your own with \emph{das Kapitel}
  \item \emph{Recall:} say the German for the plot, then the German for the character, then say \emph{das Kapitel} again
\end{itemize}

\begin{tcolorbox}[breakable,colback=teal!4,colframe=teal!35!black,title={Before you move on}]
What does das Kapitel mean? (\textquotedblleft{}The chapter\textquotedblright{}.) Say it once more.
\end{tcolorbox}
\section[die Fortsetzung]{die Fortsetzung — the sequel, the continuation}

\label{lesson:GE-C316-fortsetzung}



\begin{quote}
Before the new one: say the German for the character, then the German for the chapter.
\end{quote}

\subsection*{You'll want to know: die Fortsetzung}
\textbf{die Fortsetzung} — \textquotedblleft{}the sequel, the continuation\textquotedblright{}.

\begin{quote}
Die Fortsetzung kommt nächstes Jahr ins Kino.
\end{quote}

\begin{quote}
\emph{The sequel comes to the cinema next year.}
\end{quote}

Plural: \emph{die Fortsetzungen}. \emph{Fortsetzung folgt} — to be continued.

\begin{cousinweb}
\emph{Fortsetzung} is built on \emph{fortsetzen} \textquotedblleft{}to continue\textquotedblright{}, from \emph{setzen} \textquotedblleft{}to set\textquotedblright{}, Proto-Germanic \emph{satjaną} (EtymDB).
\end{cousinweb}

\begin{grammarlens}[title={What you've built}]
One more word for telling, explaining and describing at B1.
\end{grammarlens}

\subsection*{Your turn}
\begin{itemize}
\raggedright
  \item \emph{Say it:} \emph{die Fortsetzung}
  \item \emph{Say it:} \emph{die Fortsetzung}, once more
  \item \emph{Say it:} a sentence of your own with \emph{die Fortsetzung}
  \item \emph{Recall:} say the German for the character, then the German for the chapter, then say \emph{die Fortsetzung} again
\end{itemize}

\begin{tcolorbox}[breakable,colback=teal!4,colframe=teal!35!black,title={Before you move on}]
What does die Fortsetzung mean? (\textquotedblleft{}The sequel, the continuation\textquotedblright{}.) Say it once more.
\end{tcolorbox}
